\documentclass[11pt]{article}
\usepackage[margin=1in]{geometry}
\usepackage{amsmath,amssymb,amsthm,booktabs,longtable}
\usepackage[hidelinks]{hyperref}
\usepackage{xurl}
\setlength{\emergencystretch}{2em}
\newtheorem{theorem}{Theorem}
\newtheorem{proposition}[theorem]{Proposition}
\newtheorem{lemma}[theorem]{Lemma}
\newtheorem{corollary}[theorem]{Corollary}
\newcommand{\mr}{\operatorname{mr}}
\newcommand{\br}{\operatorname{b}}
\newcommand{\bi}{\operatorname{i}}
\title{A bounded family comparison for a seventeen-crossing knot}
\author{Primary-source comparison and written deductions}
\date{17 November 2026}
\begin{document}
\maketitle

Let $K_*$ be the oriented closure of the four-braid
\[
 \beta=\sigma_1\sigma_2\sigma_3^{-2}\sigma_2
 \sigma_3^{-2}\sigma_2\sigma_1^2\sigma_2^3
 \sigma_3^{-1}\sigma_2^3 .
\]
We compare this explicitly defined knot with specific published
meridional-rank equality criteria. The comparison separates the
certificate for $K_*$ from the hypotheses of the older constructions.
It is not an exhaustive priority claim. Write $\br$ for bridge number
and $\bi$ for braid index.

\section{Separate imported knot and reflection results}

The following results are used as explicit premises.
They are not inferred from a floating-point volume or from a census name.

\begin{proposition}[Imported knot data]\label{prop:data}
The exact diagram and quotient proof in \cite[Definition
\texttt{def:knot}, Theorems \texttt{thm:sharp} and
\texttt{thm:crossing}]{new17} gives
\[
 \mr(K_*)=\br(K_*)=4,\qquad
 g(K_*)=7,\quad \bi(K_*)=4,\quad c(K_*)=17,
\]
and the normalized Alexander polynomial
\[
\begin{split}
\Delta_{K_*}(t)={}&t^{14}-4t^{13}+8t^{12}-8t^{11}
 +18t^9-38t^8+47t^7\\
 &-38t^6+18t^5-8t^3+8t^2-4t+1 .
\end{split}
\]
In particular its breadth is $14$ and its determinant is $117$.
Theorem \texttt{thm:hyperbolic} of that companion proves that the
actual knot exterior is hyperbolic. Its separate Theorem
\texttt{thm:nongraph} proves that the actual double branched cover
$\Sigma_2(K_*)$ is not a graph manifold.
\end{proposition}

\begin{proposition}[Imported reflection obstruction]\label{prop:reflection}
Theorem \texttt{thm:reflection} and Corollary
\texttt{cor:coxeter} of \cite{universal} exclude four independent
meridian normal directions in a symmetric-form reflection
representation over characteristic zero; singular symmetric forms
are included. Consequently $K_*$ has no Coxeter reflection-class
quotient of rank at least four, finite or infinite.
\end{proposition}

The crossing-number deduction in Proposition~\ref{prop:data}
uses Yamada's theorem that the minimum Seifert-circle number of an
oriented link equals its braid index \cite[Theorem~3, p.~347]{yamada}.
Indeed, Alexander breadth gives $g\ge7$, and the displayed braid
surface gives $g\le(17-4+1)/2=7$. Bridge number four and this
four-braid give $\bi=4$. Every diagram $D$ therefore has at least
four Seifert circles, so
$c(D)=2g_D+s(D)-1\ge14+4-1=17$.
Here $g_D$ is the genus of its canonical Seifert surface.
Yamada's notation $b$ denotes braid index, rather than our $\br$.

For clarity, the closed-cover conclusion has its own volume
premises. Nondegenerate ideal shapes solving the multiplicative
gluing equations give holonomy with representation volume equal
to the signed tetrahedron sum \cite[Remark~4.4, p.~433;
Corollary~4.14 and Remark~4.15, p.~437]{franc}.
Positive tetrahedra and complete peripheral holonomy are not required.
In the actual filling certificate, the filled-slope holonomy $A$
commutes with a loxodromic peripheral holonomy $B$. The latter is
certified by a logarithmic multiplier with nonzero real part.
After conjugating $B(z)=\lambda z$, $|\lambda|\ne1$,
commutativity forces $A$ to fix $0,\infty$ individually.
The certified multiplier $1$ then gives $A=I$.
Thus the representation factors through the closed filling.
Its \emph{real} volume is unchanged by filling
\cite[Proposition~4.17, p.~437]{franc}; there is no additional real
core-length term. A certified nonzero volume contradicts
\cite[Lemma~4.10, p.~11 of the retained preprint]{lty}, which
states that every $\mathrm{PSL}_2(\mathbb C)$ representation of a
closed orientable graph manifold has volume zero.
This route uses neither the essential-edge hypothesis of that
paper's different Lemma~4.6 nor closed hyperbolicity or faithfulness.

\section{Nonalternation and positivity}

\begin{proposition}\label{prop:nonalternating}
The knot $K_*$ is nonalternating.
\end{proposition}
\begin{proof}
The polynomial in Proposition~\ref{prop:data} has zero coefficients
of $t^{10}$ and $t^4$ between nonzero extreme coefficients.
An alternating knot has no such zero Alexander coefficient
\cite[Theorem~1.1, p.~181]{murasugi}.
\end{proof}

\begin{proposition}\label{prop:rasmussen}
The knot $K_*$ is homogeneous and has Rasmussen invariant $s(K_*)=6$.
Neither it nor its mirror is positive or strongly quasipositive.
In particular it has no positive or negative braid presentation.
\end{proposition}
\begin{proof}
Each Artin generator in $\beta$ occurs with a single sign:
$\sigma_1,\sigma_2$ positively and $\sigma_3$ negatively.
The closure diagram is homogeneous
\cite[Lemma~2.2, p.~5]{abe}. Its writhe is $12-5=7$ and its
Seifert-circle number is $4$. The positive Seifert subgraph has
components $\{1,2,3\}$ and $\{4\}$; the negative subgraph has
components $\{1\},\{2\},\{3,4\}$.
The number $O_+$ of components after smoothing all negative
crossings is therefore $2$. Abe's formula
\cite[Theorem~1.2, p.~2]{abe} gives
\[
 s(K_*)=w-O+2O_+-1=7-4+4-1=6 .
\]
For the mirror, $w=-7$ and $O_+=3$, giving $s=-6$.
Both knots are homogeneous of genus seven.
For a homogeneous positive or strongly quasipositive knot,
Abe's Theorem~1.3 requires $s/2=g$; it fails for $K_*$ and its mirror.
\end{proof}

This proposition does not exclude arbitrary quasipositivity.
Proposition~\ref{prop:nonalternating} also prevents using an
alternating-diagram crosscap formula on the displayed knot.

\section{Constructions detected by Coxeter quotients}

\begin{proposition}\label{prop:coxeterfamilies}
The knot $K_*$ is outside the following precise quotient-producing
constructions:
\begin{enumerate}
\item the standard twisted-link and even-weight bipartite
arborescent theorems of Baader--Blair--Kjuchukova;
\item fishnets satisfying the nonunit even-column gcd criterion,
and the further maximal-rank Coxeter constructions in the
\emph{Other fishnets} section of Blair--Kjuchukova--Pfaff;
\item the classical integer-box families $\mathcal K_{p,4}$
of Joseph--Pongtanapaisan, including rational tangles assembled
from the permitted boxes with the stated orientation rules.
\end{enumerate}
\end{proposition}
\begin{proof}
The first two assertions follow from the rank-four Coxeter quotient
that the respective constructions supply at bridge number four:
see \cite[Theorems~1--2 and Section~4]{bbk} and
\cite[Proposition~1.4 and Sections~2.1--2.2]{fishnet}.
This contradicts Proposition~\ref{prop:reflection}.

For the third, we give the comparison rather than attributing a
Coxeter theorem to the source. In
\cite[Construction~3.5, Definition~3.6 and Corollary~3.12]{jp},
put $d=2p-1$. Replace the four standard step-$p$ cycle labels by
the standard generators of the rank-four Coxeter group with
adjacent labels $d$ and all nonadjacent labels $2$.
Equal and disjoint cycle labels become equal or commuting
reflections. For adjacent reflections $a,b$, put $q=ab$.
The two-strand involutory braid action satisfies
\[
 T(a,b)=(aba,a),\qquad
 T^k(a,b)=(q^ka,q^{k-1}a),\qquad T^d(a,b)=(a,b).
\]
Thus the allowed multiples of $d$ on coherently oriented strands
and even multiples of $d$ on oppositely oriented strands preserve the
boundary labels in this Coxeter group as well.
The four original standard labels still occur, so the labeling
is onto the rank-four Coxeter group.
The same boundary-label argument applies to the assembled
rational tangles of \cite[Remark~3.11 and Figure~5]{jp}.
Proposition~\ref{prop:reflection} again gives a contradiction.
\end{proof}

The proposition does not reject arbitrary rational replacements
of twisted diagrams or arbitrary loose fishnet diagrams.
The fishnet theorem that adds an unknot adds a link component;
it does not prove equality for the original sublink.
Also $\Omega(117)=3$, so a determinant obstruction requiring
$\Omega\le2$ would not apply here.

\section{The symmetric augmentation comparison}

\begin{lemma}\label{lem:free}
Suppose an $n$-braid $B$ representing $K_*$ has a map
$\epsilon:\mathcal A_n^{\mathrm{ab}}\to\mathbb C$ with
$\epsilon(\Phi_B^L)=\Delta(B)$ in Cornwell's conventions.
Here $\mathcal A_n^{\mathrm{ab}}$ imposes both commutativity
and $a_{ij}=a_{ji}$. Then such a map is impossible.
\end{lemma}
\begin{proof}
We use the uniform transport in
\cite[Lemma \texttt{lem:transport}]{symmetric}.
On the \emph{free} space with basis $e_1,\ldots,e_n$, put
$G(e_i,e_i)=2$ and $G(e_i,e_j)=\epsilon(a_{ij})$.
The reflections
$R_{e_i}(x)=x-G(e_i,x)e_i$ have independent displacement lines.
At a positive crossing the roots transform by
$(u,v)\mapsto(R_uv,u)$, with coefficient block
$\left(\begin{smallmatrix}-G(u,v)&1\\1&0\end{smallmatrix}\right)$.
These are the $\Phi^L$ blocks, and the Chain Rule gives final
coefficient matrix $\epsilon(\Phi_B^L)=\Delta(B)$.
Its diagonal signs close the reflections, producing a
knot-group representation on this free space.
No Gram determinant has been inverted.

The group has four generating meridians.
The span $W$ of their four displacement lines is invariant under
those reflections, and hence contains every group displacement
image. In particular it contains the $n$ independent initial
strand lines, so $n\le4$. But $\bi(K_*)=4$ gives $n\ge4$.
Thus $n=4$ and the four generating meridians also have independent
normal directions, contrary to Proposition~\ref{prop:reflection}.
\end{proof}

\begin{corollary}\label{cor:symmetric}
The knot $K_*$ has no symmetric maximal-rank augmentation of
a braid presentation in the sense used by Cornwell's
Theorem~5.4. In particular it is not an output of that
rank-raising theorem.
\end{corollary}
\begin{proof}
Cornwell's Theorem~1.10 and its proof give
$\epsilon(\Phi_B^L)=\Delta(B)$ for a rank-$n$ augmentation
of an $n$-braid. The symmetric hypothesis of
\cite[Theorem~5.4, p.~20]{cornwell} is exactly descent to
$\mathcal A_n^{\mathrm{ab}}$, so Lemma~\ref{lem:free} applies.
Theorem~5.4 preserves this symmetry in its rank-raising output.
\end{proof}

This proof does not assert that a generic augmentation retains
its rank at meridian parameter $-1$. The specialized Gram matrix
can be singular; the root directions in the free realization
are still independent. It gives no obstruction to
\emph{nonsymmetric} augmentations and no unrestricted
augmentation-rank value for $K_*$.

\section{Other bounded exclusions and remaining scope}

\begin{corollary}\label{cor:otherfamilies}
Under the separate knot-data imports in Proposition~\ref{prop:data},
$K_*$ is outside the bridge-at-most-three and torus-knot families,
the graph-manifold-exterior and nontrivial satellite output classes,
the Montesinos/arborescent graph-double-cover class, and the
bridge-four graph-cover criterion of Boileau--Jang--Weidmann.
It also fails the exact nonunit plat/matrix-gcd sharp-four
criterion surveyed by Moriah. The satellite and graph-cover
exclusions use the actual exterior and double cover, respectively;
the cable and matrix-gcd failures also follow independently from
the stated bridge and Alexander data.
\end{corollary}
\begin{proof}
Proposition~\ref{prop:data} excludes bridge at most three,
torus knots, graph-manifold exteriors and all nontrivial satellite
outputs. Its separate nongraph-cover assertion excludes the
Montesinos/arborescent graph-cover setting
\cite[Introduction, pp.~vi--vii]{bs} and the bridge-four
graph-double-cover theorem \cite[Corollary~1.3]{bjw}.
These two uses of topology should not be conflated:
a hyperbolic exterior by itself does not exclude a graph
\emph{double cover}.

There is also an independent cable check. For a genuine
winding-$p$ cable, $p\ge2$, on a nontrivial companion $J$, the satellite
bridge inequality gives
$4=\br(K_*)\ge p\br(J)\ge2p$, so $p=2$
\cite[Theorem~1, p.~4]{schultens}. The satellite Alexander formula
\cite[Theorem~2 and its Seifert corollary, p.~222]{lm}
then gives $|q|=\det(K_*)=117$ for its $(2,q)$ torus pattern.
That factor alone has Alexander breadth $116$, exceeding $14$.
An outer braid-pattern satellite must similarly have two strands,
and a two-strand braid pattern is a cable.

The exact nonunit plat/matrix-gcd criterion surveyed by Moriah
\cite[Theorem~4.20]{moriah} also cannot give a sharp-four
certificate here. Its mod-$p$ double-cover homology consequence
requires $p^3\mid\det(K_*)$ for some prime $p$; the written
deduction is \cite[Corollary \texttt{cor:plat}]{gcd}.
But $117=3^2\cdot13$. This excludes that exact criterion,
rather than every geometric generalized-Montesinos description.
\end{proof}

The broader graph-knot closure results with a meridionally
tame base \cite[Theorem~1 and Corollary~1]{bdjw}
do not automatically classify every possible hyperbolic base.
Nor does exclusion of Coxeter witnesses
exclude every representation or augmentation method.
No direct named published equality certificate was identified
in this bounded audit, but a negative search does not prove
worldwide priority. A census name is not assigned without an
exact identification. The retained source editions, pins,
family scopes and explicit remaining limitations are recorded
in \path{NEW_SEVENTEEN_LITERATURE_AUDIT.json}.

\begin{thebibliography}{99}\small
\bibitem{new17} \emph{An explicit seventeen-crossing four-bridge knot},
companion \path{NEW_SEVENTEEN_CROSSING_KNOT.tex},
Definition \texttt{def:knot}, Theorems
\texttt{thm:sharp}, \texttt{thm:crossing},
\texttt{thm:hyperbolic}, \texttt{thm:nongraph}.
\bibitem{universal} \emph{Universal reflection obstruction},
companion \path{UNIVERSAL_17_REFLECTION_OBSTRUCTION.tex},
Theorem \texttt{thm:reflection}, Corollary \texttt{cor:coxeter}.
\bibitem{yamada} S.~Yamada, \emph{The minimal number of Seifert
circles equals the braid index of a link}, Invent.\ Math.\
89 (1987), 347--356, doi:10.1007/BF01389082.
\bibitem{franc} S.~Francaviglia, \emph{Hyperbolic volume of
representations of fundamental groups of cusped 3-manifolds},
IMRN 2004(9), 425--459, doi:10.1155/S1073792804131619.
\bibitem{lty} F.~Luo, S.~Tillmann and T.~Yang,
\emph{Thurston's spinning construction and solutions to the
hyperbolic gluing equations for closed hyperbolic 3-manifolds},
arXiv:1004.2992v1 (2010), Lemma~4.10.
\bibitem{murasugi} K.~Murasugi, \emph{On the Alexander polynomial
of the alternating knot}, Osaka Math.\ J.\ 10 (1958), 181--189.
\bibitem{abe} T.~Abe, \emph{The Rasmussen invariant of a homogeneous
knot}, Proc.\ AMS 139 (2011), 2647--2656; retained arXiv:1003.5392v1.
\bibitem{bbk} S.~Baader, R.~Blair and A.~Kjuchukova,
\emph{Coxeter groups and meridional rank of links},
Math.\ Ann.\ 379 (2021), 1533--1551.
\bibitem{fishnet} R.~Blair, A.~Kjuchukova and C.~Pfaff,
\emph{Adding a suitable unknot to any link equates bridge number
and meridional rank}, doi:10.1017/S030500412610214X (2026);
retained arXiv:2411.10642 author version.
\bibitem{jp} J.~Joseph and P.~Pongtanapaisan, \emph{Meridional
rank and bridge number of knotted 2-spheres},
Canad.\ J.\ Math.\ 77 (2025), 282--299; retained arXiv:2111.09233v2,
under the title \emph{Bridge number and meridional rank of knotted
surfaces}.
\bibitem{cornwell} C.~R.~Cornwell, \emph{KCH representations,
augmentations, and $A$-polynomials}, J.\ Symplectic Geom.\ 15 (2017),
983--1017; retained arXiv:1310.7526v3.
\bibitem{symmetric} \emph{Symmetric maximal-braid augmentations
and free reflection directions}, companion
\path{SYMMETRIC_AUGMENTATION_OBSTRUCTION.tex},
Lemma \texttt{lem:transport}.
\bibitem{bjw} M.~Boileau, Y.~Jang and R.~Weidmann,
\emph{Meridional rank and bridge number for a class of links},
Pacific J.\ Math.\ 292 (2018), 61--75.
\bibitem{bs} F.~Bonahon and L.~C.~Siebenmann,
\emph{New geometric splittings of classical knots and the
classification and symmetries of arborescent knots},
author monograph, version 26 January 2016.
\bibitem{bdjw} M.~Boileau, E.~Dutra, Y.~Jang and R.~Weidmann,
\emph{Meridional rank of knots whose exterior is a graph manifold},
arXiv:1608.01570.
\bibitem{schultens} J.~Schultens, \emph{Additivity of bridge numbers
of knots}, Math.\ Proc.\ Camb.\ Phil.\ Soc.\ 135 (2003), 539--544.
\bibitem{lm} C.~Livingston and P.~Melvin, \emph{Abelian invariants
of satellite knots}, Geometry and Topology, LNM 1167 (1985), 217--227.
\bibitem{moriah} Y.~Moriah, \emph{Heegaard splittings of knot
exteriors}, Geom.\ Topol.\ Monogr.\ 12 (2007), 191--232,
Theorem~4.20.
\bibitem{gcd} \emph{Fishnet scope audit}, companion
\path{FISHNET_SCOPE_AUDIT.tex}, Corollary \texttt{cor:plat}.
\end{thebibliography}
\end{document}
